\chapter{Ecuaciones de diferencias finitas}

\section{Introducci\'{o}n}

En los sistemas lineales de tiempo discreto invariantes en el tiempo la entrada y la salida est\'{a}n 
relacionadas por 
\begin{enumerate}
   \item  una ecuaci\'{o}n de diferencias 
      \begin{equation}
           \sum_{p=0}^{N}a_{p} \, y\left[ n-p\right] =\sum_{q=0}^{M}b_{q} \, x\left[ n-q\right]. 
           \label{eqdiferencia}
      \end{equation}

   \item   la suma de convoluci\'{o}n
        \begin{equation}
                y[n] = \sum_{k=-\infty}^{\infty} x[k] \, h[n - k].
        \end{equation}
        
\end{enumerate}
Ambas formas deben ser equivalentes. Justamente esta equivalencia da pie a que podamos usar la
transformada $Z$ para resolver ecuaciones de diferencias. Empezaremos estudiando la equivalencia entre 
ambas formas de c\'{a}lculo para entradas de la forma $x[n] = z^{n}$, y despu\'{e}s generalizaremos 
considerando entradas para las que existan las transformadas $Z$ correspondientes.

\section{Ecuaci\'{o}n de diferencias con entrada exponencial compleja}

Consideremos una ecuaci\'{o}n de diferencias de la forma 
\begin{equation}
   \sum_{p=0}^{N}a_{p} \, y\left[ n-p\right] =\sum_{q=0}^{M}b_{q} \, x\left[ n-q\right].
\end{equation}
Buscaremos la soluci\'{o}n para entradas de la forma 
\begin{equation}
x[n]=z^n ,
\end{equation}
donde $z$ es una variable compleja.

\subsection{Soluci\'{o}n de una ecuaci\'{o}n de diferencias homog\'{e}nea}

\begin{enumerate}

\item Denominaremos $y_{h}\left[ n\right] $ a la soluci\'{o}n de la ecuaci\'{o}n homog\'{e}nea
\begin{equation}
        \sum_{p=0}^{N}a_{p} \, y_{h}\left[ n-p\right] =0.
\end{equation}

\item Supongamos que la soluci\'{o}n tiene la forma 
\begin{equation}
        y_{h}\left[ n\right] = Y \, \lambda^n,
\end{equation}
donde $Y$ es una variable real. Observemos que 
\begin{equation}
    y_{h}\left[ n-p\right] = Y \, \lambda^{n-p},
\end{equation}
para todos los valores de $p$ en el intervalo $[0,N]$.

\item Sustituyendo en la ecuaci\'{o}n diferencial homog\'{e}nea obtenemos 
\begin{equation}
    \sum_{p=0}^{N}a_{p} \, Y \, \lambda^{n-p}  =\left( \sum_{p=0}^{N}a_{p} \, \lambda^{-p} \right) \, Y \, \lambda^{n}=0.
\end{equation}
No estamos interesados en la soluci\'{o}n trivial $Y=0$. Como la se\~{n}al $\lambda^{n}$ no se anula para 
ning\'{u}n valor de $n$, tiene que ser 
\begin{equation}
  \left( \sum_{p=0}^{N}a_{p} \, \lambda ^{-p}\right) =
  \lambda^{-N} \, \left( \sum_{p=0}^{N}a_{p} \, \lambda ^{N-p}\right) = 0.
\end{equation}

\item Definimos entonces al polinomio caracter\'{\i}stico $p(\lambda)$
\begin{equation}
  p(\lambda) = \left( \sum_{p=0}^{N}a_{p} \, \lambda ^{N-p}\right) = 
  a_{0} \, \prod_{q=1}^{N} \left( \lambda - \lambda_{q} \right) = 0.
\end{equation}
Las ra\'{\i}ces $\lambda_{q}$ del polinomio caracter\'{\i}stico son llamadas 
frecuencias, modos naturales o polos.

\item Si las frecuencias o modos naturales son distintos 
\begin{equation}
        \left( i\neq j\longleftrightarrow \lambda _i\neq \lambda _{j}\right),
\end{equation}
entonces la soluci\'{o}n m\'{a}s general tiene la forma 
\begin{equation}
        y_{h}\left[ n\right] =\sum_{k=1}^{N}Y_{k} \, \lambda _{k}^n.
\end{equation}

\item Si hay $R<N$ frecuencias o modos naturales distintos la soluci\'{o}n homog\'{e}nea tiene la forma 
\begin{equation}
        y_{h}\left[ n\right] =\sum_{p=1}^{R}\left(\sum_{k=0}^{S_{p}-1}Y_{p,k} \, n^{k}\right) \, \lambda _{p}^n,
\end{equation}
donde $S_{p}$ indica la multiplicidad de $\lambda _{p}$ y $\sum_{p=1}^{R}S_{p}=N$.

\end{enumerate}

\subsubsection{Ejemplo}

\begin{enumerate}

\item Consideremos la ecuaci\'{o}n de diferencias
\begin{equation}
  -v_{o}\left[ n+1\right] +\left( 2+c\right) \, v_{o}\left[ n\right] -v_{o} \, \left[n-1\right] =
     c \, v_{i}\left[ n\right].
\end{equation}

\item La ecuaci\'{o}n homog\'{e}nea asociada es
\begin{equation}
  -v_{o} \, \left[ n+1\right] +\left( 2+c\right) \, v_{o}\left[ n\right] -v_{o}\left[n-1\right] = 0.
\end{equation}

\item Suponiendo que
\begin{equation}
  \begin{array}{ccc}
    v_{o}[n+1] = V \, \lambda^{n+1}, & v_{o}[n] = V \, \lambda^{n}, & v_{o}[n-1] = V \, \lambda^{n-1}, 
  \end{array}
\end{equation}
reemplazamos en la ecuaci\'{o}n homog\'{e}nea
\begin{equation}
  -V \, \lambda^{n+1} + (2 + c) \, V \, \lambda^{n} - V \, \lambda^{n-1} = 0. \label{ecdifhomuno}
\end{equation}

\item Extrayendo el factor com\'{u}n $V \, \lambda^{n-1}$ obtenemos el polinomio caracter\'{\i}stico
\begin{equation}
  V \, \lambda^{n-1} \, (-\lambda^{2} + (2 + c) \, \lambda^{1} - 1) = 0. \label{ecdifhomdos}
\end{equation}

\item Las frecuencias naturales son 
\begin{equation}
        \lambda_{1,2}=\left( 1+\frac{c}{2}\right) \pm \sqrt{\left( 1+\frac{c}{2}\right)^{2}-1}.
\end{equation}

\item Observemos que
\begin{equation}
        \left( 
                \left( 1+\frac{c}{2}\right) +
                \sqrt{\left( 1+\frac{c}{2}\right)^{2}-1}
        \right) 
        \left( \left( 1+\frac{c}{2}\right) -\sqrt{\left( 1+\frac{c}{2}\right) ^{2}-1}
        \right) =1.
\end{equation}
y por lo tanto 
\begin{equation}
        \lambda _{1}=\frac{1}{\lambda _{2}}.
\end{equation}

\item Entonces la soluci\'{o}n homog\'{e}nea es 
\begin{equation}
        v_{0}\left[ n\right] = V_{1} \, \lambda _{1}^n+ V_{2} \, \lambda _{1}^{-n}.
\end{equation}

\item Observe las ecs. (\ref{ecdifhomuno}) y (\ref{ecdifhomdos}). 
\textquestiondown Porqu\'{e} el factor es $V \, \lambda^{n-1}$ en lugar de $V \, \lambda^{n-2}$, 
dado que $N=2$?
\end{enumerate}

\subsection{Soluci\'{o}n particular de una ecuaci\'{o}n de diferencias para una entrada $z^{n}$}

\subsubsection{Donde la entrada $z^{n}$ no es una frecuencia natural}

\begin{enumerate}

\item Consideremos la ecuaci\'{o}n de diferencias
\begin{equation}
  \sum_{p=0}^{N}a_{p} \, y\left[ n-p\right] =
  \sum_{q=0}^{M}b_{q} \, x\left[ n-q\right]. \label{ecdiflinealnonat}
\end{equation}

\item Supongamos ahora que la entrada es una exponencial
\begin{equation}
  x\left[ n\right] = z^{n},
\end{equation}
y que la salida tiene la misma forma
\begin{equation}
  y\left[ n\right] = Y \, z^{n},
\end{equation}
donde $z$ no es igual a ninguna frecuencia natural. 
Reemplazando en la ecuaci\'{o}n de diferencias resulta 
\begin{equation}
  \sum_{p=0}^{N}a_{p} \, Y \, z^{n-p}=\sum_{q=0}^{M}b_{q} \, z^{n-q}.
\end{equation}

\item Extrayendo factores comunes obtenemos 
\begin{equation}
  \left( \sum_{p=0}^{N}a_{p} \, z^{-p}\right) \, Y \, z^{n}=
  \left(\sum_{q=0}^{M}b_{q} \, z^{-q}\right) \, z^{n}. \label{ecdifpart}
\end{equation}
Luego de dividir por $z^n$, podemos despejar $Y$
\begin{equation}
  Y=\frac{\left( \sum_{q=0}^{M}b_{q} \, z^{-q}\right) }{\left(\sum_{p=0}^{N}a_{p} \, z^{-p}\right) }.
\end{equation}

\item Por lo tanto, una soluci\'{o}n particular de la ecuaci\'{o}n (\ref{ecdiflinealnonat}) 
para una entrada exponencial $z^{n}$ es
\begin{equation}
  y_{p}[n] =
    \frac{\left( \sum_{q=0}^{M}b_{q} \, z^{-q}\right) }{\left(\sum_{p=0}^{N}a_{p} \, z^{-p}\right) } \, z^{n}.
    \label{eqhzorig}
\end{equation}

\end{enumerate}

\subsubsection{Donde la entrada $z^{n}$ es una frecuencia natural}

\begin{enumerate}

\item Consideremos la ecuaci\'{o}n de diferencias
\begin{equation}
  \sum_{p=0}^{N}a_{p} \, y\left[ n-p\right] =
  \sum_{q=0}^{M}b_{q} \, x\left[ n-q\right]. \label{ecdiflinealnat}
\end{equation}

\item Supongamos ahora que la entrada es una frecuencia natural
\begin{equation}
  x\left[ n\right] = \lambda^{n}.
\end{equation}
Por lo tanto debe cumplirse (ver ec. \ref{ecdifpart})
\begin{equation}
  \left( \sum_{p=0}^{N}a_{p} \, \lambda^{-p}\right) = 0,
\end{equation}
anul\'{a}ndose as\'{\i} el \textbf{denominador} de la ec. (\ref{eqhzorig}).

\item Este es un fen\'{o}meno denominado resonancia. Se produce cuando las entradas coinciden con las
frecuencias naturales del sistema, y provocan una amplificaci\'{o}n de la salida. Esta ha sido la causa
de la destrucci\'{o}n de puentes y aviones. Busque referencias en Internet sobre el puente sobre el 
r\'{\i}o Tacoma, o Tacoma River, o Tacoma Bridge.

\item En este caso, la respuesta forzada del sistema es
\begin{equation}
  y[n] = \left(n^{M} + \sum_{k=0}^{M-1} c_{k} \, n^{k}\right) \lambda^{n}, \label{resonanciadiscreta}
\end{equation}
donde $M$ es la multiplicidad de la frecuencia natural $\lambda$, y los coeficientes $c_{k}$ se
determinan mediante condiciones iniciales.

\item Por ejemplo, consideremos la ecuaci\'{o}n de diferencias
\begin{equation}
  y[n] - \lambda y[n-1]= \lambda^{n}. 
\end{equation}
Podemos comprobar que sustituyendo
\begin{eqnarray}
  y[n] &=&  (n + c) \, \lambda^{n}, \label{resonanciadiscretarespuesta} \\
  y[n-1] &=&  ((n - 1) + c) \, \lambda^{n-1},
\end{eqnarray}
en la ec. (\ref{resonanciadiscreta}), obtenemos
\begin{equation}
  (n + c) \, \lambda^{n} - \lambda \, ((n - 1) + c) \, \lambda^{n-1} = \lambda^{n}.
\end{equation}
Reagrupando y simplificando, resulta
\begin{equation}
   n \, \lambda^{n} - n \, \lambda^{n} + c \, \lambda^{n} - c \, \lambda^{n} + \lambda^{n} = \lambda^{n}.
\end{equation}
Por lo tanto la respuesta forzada es de la forma (\ref{resonanciadiscreta}).

\end{enumerate}

\subsection{Soluci\'{o}n general de la ecuaci\'{o}n de diferencias para una entrada $z^{n}$}

\subsubsection{Donde $z^{n}$ no es una frecuencia natural}

La soluci\'{o}n general buscada de la ecuaci\'{o}n de diferencias
\begin{equation}
        \sum_{p=0}^{N}a_{p}y\left[ n-p\right] =\sum_{q=0}^{M}b_{q}x\left[ n-q\right].
\end{equation}
es de la forma
\begin{equation}
   y[n] = y_{h} + y_{p}[n] = 
     \left(\sum_{k=1}^{N}Y_{k}\lambda _{k}^n \right) +
    \frac{\left( \sum_{q=0}^{M}b_{q}z^{-q}\right) }{\left(\sum_{p=0}^{N}a_{p}z^{-p}\right) } z^{n}.
\end{equation}

\subsubsection{Donde la entrada $z^{n}$ es una frecuencia natural}

\section{Definici\'{o}n de la funci\'{o}n de transferencia para sistemas de tiempo discreto}

Definamos entonces 
\begin{equation}
  H\left( z\right) =
  \frac{\left( \sum_{q=0}^{M}b_{q}z^{-q}\right) }{\left(\sum_{p=0}^{N}a_{p}z^{-p}\right) },
  \label{eqhzdef}
\end{equation}
como la funci\'{o}n de transferencia asociado a la ecuaci\'{o}n de diferencias
\begin{equation}
  \sum_{p=0}^{N}a_{p}y\left[ n-p\right] =
  \sum_{q=0}^{M}b_{q}x\left[ n-q\right].  
\end{equation}

\subsection{Ejemplo uno}

\begin{enumerate}

\item La ecuaci\'{o}n de diferencias
\begin{equation}
  y[n] - y[n-1] = \tau x[n],
\end{equation}
tiene una funci\'{o}n de transferencia asociada
\begin{equation}
  H(z) = \frac{\tau}{1-z^{-1}}.
\end{equation}

\item La ecuaci\'{o}n de diferencias
\begin{equation}
  y[n] - y[n-2] = 2 \tau x[n-1],
\end{equation}
tiene una funci\'{o}n de transferencia asociada
\begin{equation}
  H(z) = \frac{2 \tau z^{-1}}{1-z^{-2}}.
\end{equation}

\end{enumerate}

\subsection{Ejemplo dos}

Veremos la reconstrucci\'{o}n de una ecuaci\'{o}n de diferencias a partir de
una funci\'{o}n de transferencia. Consideremos 
\begin{equation}
        H\left( z\right) =\frac{z}{z+1}.
\end{equation}
\textquestiondown Cu\'{a}l es la ecuaci\'{o}n de diferencias correspondiente?
Hagamos 
\begin{equation}
        H(z)\left( z+1\right) = z,
\end{equation}
y multipliquemos ambos lados por $z^n$
\begin{equation}
        H(z)z^n\left( z+1\right) =z^{n}z,
\end{equation}
de donde se puede obtener la ecuaci\'{o}n diferencia 
\begin{equation}
        y\left[ n+1\right] +y\left[ n\right] = z^{n+1} .
\end{equation}

\subsection{Pregunta}

Compare las funciones de transferencia con sus respectivas ecuaciones de diferencias.
\textquestiondown Encuentra la relaci\'{o}n existente entre los exponentes de $z$ en las funciones de 
transferencia, los \'{\i}ndices de los coeficientes y los desplazamientos en el tiempo de la entrada
y la salida?

\section{Polos y ceros de la funci\'{o}n de transferencia}

\begin{enumerate}

\item La funci\'{o}n de transferencia puede expresarse como 
\begin{equation}
H\left( z\right) =
        \frac{\left( \sum_{q=0}^{M}b_{q}z^{q}\right) }{\left(\sum_{p=0}^{N}a_{p}z^{p}\right) }=
        K\frac{\left(\Pi _{q=1}^{M} z-c_{q} \right)}
         {\left(\Pi _{p=1}^{N} z-d_{p} \right)},
\end{equation}
donde $K=\frac{b_{M}}{a_{N}}$.
Nos referiremos a $c_{q}$ como los ceros, 
a $d_{p}$ como los polos, y a $K$ como la ganancia.

\item La posici\'{o}n de los ceros y polos en el plano complejo determinan el tipo de din\'{a}mica 
del sistema. El factor de ganancia $K$ s\'{o}lo determina la amplitud de la respuesta.

\item La funci\'{o}n de transferencia $H\left( z\right) $ relacionada con una ecuaci\'{o}n de 
diferencia contiene la misma informaci\'{o}n que esta. Ambas est\'{a}n definidas por $N+M+1$ 
n\'{u}meros: $N$ polos, $M$ ceros y una constante o ganancia.

\end{enumerate}

\subsection{Soluci\'{o}n general de la ecuaci\'{o}n de diferencias con entrada exponencial compleja}

\begin{enumerate}

\item La soluci\'{o}n general de la ecuaci\'{o}n de diferencias 
\begin{equation}
        \sum_{p=0}^{N}a_{p}y\left[ n-p\right] =\sum_{q=0}^{M}b_{q}x\left[ n-q\right],
\end{equation}
con entrada exponencial compleja $x[n] = z^{n}$ es
\begin{equation}
  y[n] = H(z) z^{n} + y_{h}[n],
\end{equation}
donde
\begin{equation}
     H(z) = \frac{\left( \sum_{q=0}^{M}b_{q}z^{-q}\right) }{\left(\sum_{p=0}^{N}a_{p}z^{-p}\right) },
\end{equation}
y $y_{h}[n]$ satisface la ecuaci\'{o}n homog\'{e}nea
\begin{equation}
        \sum_{p=0}^{N}a_{p}y_{h}\left[ n-p\right] = 0.
\end{equation}

\item Para determinar por completo la soluci\'{o}n total necesitamos determinar los 
$N$ coeficientes $A_{1},$ ..., $A_{N}.$ que forman parte de la soluci\'{o}n homog\'{e}nea $y_{h}[n]$.
Estos pueden ser determinados a partir de $N$ condiciones iniciales que deben ser especificadas. 
Se obtiene as\'{\i} un conjunto de $N$ ecuaciones algebraicas cuyas inc\'{o}gnitas
son los $N$ coeficientes $A_{1},$ ..., $A_{N}$.

\end{enumerate}

\section{Respuesta a exponenciales complejas}

\begin{enumerate}

\item Sea el sistema definido por la suma de convoluci\'{o}n, 
\begin{equation}
        y\left[ n\right] =
                \sum_{k=-\infty }^{\infty }x\left[ k\right] h\left[ n-k\right] =
                \sum_{k=-\infty }^{\infty }h\left[ k\right] x\left[ n-k\right] ,
\end{equation}
donde $h\left[ n\right] $ es la respuesta al impulso, $x\left[ n\right] $ es la entrada e 
$y\left[ n\right] $ es la salida correspondiente.

\item Sup\'{o}ngamos una entrada exponencial compleja 
\begin{equation}
        x\left[ n\right] =z^{n}.
\end{equation}
Por lo tanto, la suma de convoluci\'{o}n correspondiente es 
\begin{equation}
        y\left[ n\right] =
                \sum_{k=-\infty }^{\infty }h\left[ k\right] x\left[ n-k\right] =
                \sum_{k=-\infty }^{\infty }h\left[ k\right] z^{n-k},
\end{equation}
que puede expresarse como 
\begin{equation}
y\left[ n\right] =\left( \sum_{k=-\infty }^{\infty }h\left[ k\right]z^{-k}\right) z^n.
\end{equation}

\item Vemos que la funci\'{o}n definida por
\begin{equation}
        H\left( z\right) =\left( \sum_{k=-\infty }^{\infty }h\left[ k\right]z^{-k}\right),
\end{equation}
es la transformada Z bilateral de la respuesta al impulso del sistema.

\item \textbf{Muy importante: La transformada Z de la respuesta al impulso del sistema es la 
funci\'{o}n de transferencia del sistema.}

\end{enumerate}

\section{Soluci\'{o}n general de la ecuaci\'{o}n de diferencias}

\subsection{Introducci\'{o}n}

\begin{enumerate}

\item Hemos visto que podemos obtener la funci\'{o}n de transferencia de un sistema 
(la transformada $Z$ de la respuesta al impulso) suponiendo que la entrada es $x[n] = z^{n}$. 
En este caso la salida es $y[n] = H(z) z^{n}$.

\item Adem\'{a}s, una soluci\'{o}n particular de la ecuaci\'{o}n de diferencias 
est\'{a} dada por la convoluci\'{o}n de la respuesta al impulso $h[n]$ con la 
entrada $x[n]$. 

\item Disponiendo de $H(z)$ es posible evitar el c\'{a}lculo de la convoluci\'{o}n, 
porque una propiedad de la transformada $Z$ es que la transformada de una 
convoluci\'{o}n de dos se\~{n}ales $h[n]$ y $x[n]$ es el producto de las 
transformadas $H(z)$ y $X(z)$. Por lo tanto podemos:
\begin{enumerate}
   \item Calcular $X(z)$.
   \item Obtener $y_{p}[n]$ tal que $Z(y_{p}[n]) = H(z) X(z)$ mediante la 
     transformada $Z$ bilateral, realizando la suposici\'{o}n correspondiente 
     sobre la regi\'{o}n de convergencia y por lo tanto sobre la causalidad o no 
     del sistema lineal descripto por la ecuaci\'{o}n de diferencias.
\end{enumerate}

\item Sumando la soluci\'{o}n particular $y_{p}[n]$ encontrada en el paso anterior a 
  la soluci\'{o}n de la ecuaci\'{o}n homog\'{e}nea $y_{h}[n]$, obtenemos finalmente 
  la soluci\'{o}n general $y[n]$.

\end{enumerate}

\subsection{Sistemas causales}

\begin{enumerate}

\item Supondremos que el sistema definido por la ecuaci\'{o}n de diferencias 
\begin{equation}
   \sum_{p=0}^{N}a_{p}y\left[ n-p\right] =\sum_{q=0}^{M}b_{q}x\left[ n-q\right]. \label{eqdtdcausal}
\end{equation}
es \textbf{causal}, o sea que su respuesta al impulso es nula para $n < 0$.

\item Comprobaremos que la soluci\'{o}n causal de la ecuaci\'{o}n de diferencias (\ref{eqdtdcausal}) es
\begin{equation}
   y[n] = \left( \sum_{k=1}^{N} c_{k} \lambda_{k}^{n} u[n] \right) +
     \left( \sum_{k=0}^{n} h[n-k] x[k] \right). \label{solgenecdif}
\end{equation}
donde $\lambda_{k}$, $k \in [1,N]$ son las $N$ ra\'{\i}ces simples del polinomio caracter\'{\i}stico
\begin{equation}
  p(\lambda) = \left( \sum_{p=0}^{N}a_{p}\lambda ^{N-p}\right) = 
  a_{N} \prod_{k=1}^{N} \left( \lambda - \lambda_{k} \right). \label{eqdtdcausalpolcar}
\end{equation}

\item Supondremos que la respuesta al impulso es una combinaci\'{o}n lineal de funciones 
$\lambda_{k}^{n} u[n]$
\begin{equation}
   h[n] = \sum_{k=1}^{N} d_{k} \lambda_{k}^{n} u[n].
\end{equation}
donde $\lambda_{k}$, $k \in [1,N]$, son ra\'{\i}ces del polinomio caracter\'{\i}stico 
(\ref{eqdtdcausalpolcar}).

\item Recordemos que hemos llegado a la conclusi\'{o}n de que la transformada $Z$ de la respuesta causal 
al impulso es la funci\'{o}n de transferencia $H(z)$ definida como
\begin{equation}
  H\left( z\right) =
  \frac{\left( \sum_{q=0}^{M}b_{q}z^{-q}\right) }{\left(\sum_{p=0}^{N}a_{p}z^{-p}\right) } = 
  \sum_{k=0 }^{\infty }h\left[ k\right]z^{-k}.
\end{equation}
Tomando la transformada $Z$ de $h[n]$ encontramos
\begin{equation}
  Z(h[n]) = \sum_{k=-\infty}^{\infty} \left(\sum_{s=1}^{N} d_{s} \lambda_{s}^{k} u[k] \right) z^{-k}.
\end{equation}
Intercambiando la posici\'{o}n de las sumatorias resulta
\begin{equation} \label{ecuazetabilateral}
  Z(h[n]) = \sum_{s=1}^{N} d_{s} \left( \sum_{k=-\infty}^{\infty} \lambda_{s}^{k} u[k] z^{-k} \right).
\end{equation}
Es posible cambiar el valor inicial de $k$ debido al factor $u[k]$,
\begin{equation} \label{ecuazetaunilateral}
  Z(h[n]) = \sum_{s=1}^{N} d_{s} \left( \sum_{k=0}^{\infty} \lambda_{s}^{k} z^{-k} \right).
\end{equation}
Por lo tanto debe ser
\begin{equation}
  H\left( z\right) = Z(h[n]) =
  \frac{\left( \sum_{q=0}^{M}b_{q}z^{-q}\right) }{\left(\sum_{p=0}^{N}a_{p}z^{-p}\right) } = 
  \sum_{s=1}^{N} d_{s} \frac{1}{1 - (\lambda_{s} z^{-1})},
\end{equation}
donde $\left \|(\lambda_{s} z^{-1}) \right\| < 1$, $s \in [1,N]$, en la regi\'{o}n de convergencia.

\item Observemos la diferencia entre las ecuaciones 
  (\ref{ecuazetabilateral}) y (\ref{ecuazetaunilateral}). Hemos supuesto que el sistema es causal, y 
por lo tanto en su respuesta al impulso aparece el factor $u[n]$ en estas ecuaciones. 
Se justifica as\'{\i} el empleo de la transformada $Z^{+}$ (unilateral).

\item Si la entrada es un impulso, y las condiciones iniciales son nulas (o sea, 
$c_{k} = 0$) obtenemos
\begin{equation}
   h[n] = \sum_{k=0}^{\infty} h[k] \delta[n-k].
\end{equation}

\item Podemos comprobar que la soluci\'{o}n $y[n]$ tiene la forma indicada en la ec. (\ref{solgenecdif}) 
por dos caminos distintos: en el dominio del tiempo y en el dominio $z$. 

\item En el dominio del tiempo, podemos escribir
\begin{equation}
   y[n-p] = 
       \left( \sum_{k=1}^{N} c_{k} \lambda_{k}^{n-p} u[n-p] \right) +
       \left( \sum_{k=0}^{n-p} h[n-k-p] x[k] \right),
\end{equation}
Reemplazando en la ecuaci\'{o}n de diferencias, resulta
\begin{equation}
   \sum_{p=0}^{N}a_{p}y\left[ n-p\right] =
   \left( \begin{array}{c}
       \sum_{p=0}^{N}a_{p}
       \left( \sum_{k=1}^{N} c_{k} \lambda_{k}^{n-p} u[n-p] \right) + \\
       \\
       \sum_{p=0}^{N}a_{p}
       \left( \sum_{k=0}^{n-p} h[n-k-p] x[k] \right)
 \end{array} \right).
\end{equation}
Observemos que podemos escribir
\begin{equation}
  \sum_{p=0}^{N}a_{p}
  \left( \sum_{k=1}^{N} c_{k} \lambda_{k}^{n-p} u[n-p] \right) =
  \sum_{k=1}^{N} c_{k} \lambda_{k}^{n} \sum_{p=0}^{N} a_{p} \lambda_{k}^{-p} u[n-p].
\end{equation}

\end{enumerate}

\subsection{Resumen: La funci\'{o}n de transferencia}

Debemos recordar que
\begin{enumerate}
\item   $H(z)$ es una funci\'{o}n racional en $z$ para ecuaciones de diferencias con coeficientes
  constantes y sin retardos.
  
\item   $H(z)$ caracteriza al sistema lineal.
  
\item   $H(z)$ define la relaci\'{o}n entre entrada y salida.
  
\item   $H(z)$ se calcula de dos formas:
  \begin{enumerate}
  \item Considerando una entrada exponencial compleja $x[n] = z^{n}$.
  \item Aplicando la transformada $Z$ unilateral con condiciones iniciales nulas a 
    la ecuaci\'{o}n diferencial.
  \end{enumerate}
\end{enumerate}

\subsection{Ejemplo}

\begin{enumerate}

\item Hemos considerado anteriormente una versi\'{o}n discretizada del sistema de tiempo continuo cuya 
ecuaci\'{o}n diferencial es 
\begin{equation}
        \frac{dv_o\left( t\right) }{dt}+\frac{1}{RC}v_o\left( t\right) =\frac{1}{RC}v_i\left( t\right) ,
\end{equation}
y su respuesta al escal\'{o}n, comenzando de $v_o\left( 0\right) =0,$ es 
\begin{equation}
        v_o\left( t\right) =\left( 1-e^{-\frac{t}{RC}}\right) u\left( t\right) .
\end{equation}
Hemos mostrado que una aproximaci\'{o}n discreta de este sistema permite obtener la ecuaci\'{o}n de 
diferencias 
\begin{equation}
        v_o\left[ n+1\right] =\left( 1-\alpha \right) v_o\left[ n\right] +\alpha v_i\left[ n\right] ,
\end{equation}
donde 
\begin{equation}
        \alpha =\frac{T}{RC}.
\end{equation}

\item Podemos obtener la funci\'{o}n de transferencia del sistema sustituyendo 
\begin{equation}
        v_i\left[ n\right] =z^n,
\end{equation}
y 
\begin{equation}
        v_o\left[ n\right] =V_o z^n,
\end{equation}
en la ecuaci\'{o}n de diferencias para obtener 
\begin{equation}
        V_oz^{n+1}=\left( 1-\alpha \right) V_oz^n+\alpha z^n,
\end{equation}
que puede simplificarse 
\begin{equation}
        V_oz=\left( 1-\alpha \right) V_o+\alpha.
\end{equation}
Por lo tanto 
\begin{equation}
        H\left( z\right) = V_o\left( z\right) = \frac{\alpha }{z-\left( 1-\alpha \right) }.
\end{equation}

\item La frecuencia natural del sistema es 
\begin{equation}
        \lambda =\left( 1-\alpha \right) ,
\end{equation}
y por lo tanto la soluci\'{o}n tiene la forma 
\begin{equation}
     v_o\left[ n\right] =\left( A\left( 1-\alpha \right) ^n+H\left( 1\right) 1^n\right) u\left[ n\right] .
\end{equation}

\item Como $H\left( 1\right) =1,$ resulta 
\begin{equation}
  v_o\left[ n\right] =\left( A\left( 1-\alpha \right) ^n+1\right) u\left[ n\right] .
\end{equation}
Finalmente, la condici\'{o}n inicial $v_o\left[ 0\right] =0$ implica que 
$A=-1$ y por lo tanto la soluci\'{o}n es 
\begin{equation}
  v_o\left[ n\right] =\left( 1-\left( 1-\alpha \right) ^n\right) u\left[ n\right] .
\end{equation}

\end{enumerate}

\section{Estabilidad de sistemas LIT}

\subsection{Sistemas con un polo}

La funci\'{o}n de transferencia correspondiente a la ecuaci\'{o}n
diferencial 
\begin{equation}
        y\left[ n+1\right] +ay\left[ n\right] =bx\left[ n\right] ,
\end{equation}
puede calcularse suponiendo que 
\begin{eqnarray}
        x\left[ n\right] &=&Xz^n, \\
        y\left[ n\right] &=&Yz^n\longleftrightarrow y\left[ n+1\right] =Yz^{n+1}.
\end{eqnarray}
Reemplazando en la ecuaci\'{o}n diferencial 
\begin{equation}
        Yz^{n+1}+aYz^n=bXz^n
\end{equation}
Simplificando y agrupando 
\begin{equation}
        \left( z+a\right) Y=bX.
\end{equation}
Despejando $\frac{Y}{X}$, resulta 
\begin{equation}
        H\left( z\right) =\frac{Y}{X}=\frac{b}{z+a},
\end{equation}
Recuerde que la soluci\'{o}n de la ecuaci\'{o}n homog\'{e}nea es 
\begin{equation}
        y\left[ n\right] =Y\left( -a\right) ^n.
\end{equation}
\begin{enumerate}
        \item   Si $\left| a\right| < 1$ entonces 
        \begin{equation}
                \begin{array}{c}
                        Lim \\ 
                        t\rightarrow \infty
                \end{array}
                y\left[ n\right] =0,
        \end{equation}
        y por lo tanto el sistema es estable. 
        
        \item   Si $\left| a\right| > 1,$ entonces 
        \begin{equation}
                \begin{array}{c}
                        Lim \\ 
                        n\rightarrow \infty
                \end{array}
                y\left[ n\right] =\infty .
        \end{equation}
        y el sistema es inestable.
        
        \item   Si $\left| a\right| = 1,$ entonces 
        \begin{equation}
                \begin{array}{c}
                        Lim \\ 
                        n\rightarrow \infty
                \end{array}
                y\left[ n\right] = 1,
        \end{equation}
        y el sistema es estable en el sentido de que una entrada acotada genera una salida limitada.
\end{enumerate}

\subsection{Sistemas con dos polos}

La funci\'{o}n de transferencia correspondiente a la ecuaci\'{o}n de
diferencias 
\begin{equation}
        y\left[ n+2\right] +ay\left[ n+1\right] +by\left[ n\right] =cx\left[ n\right],
\end{equation}
puede calcularse suponiendo que 
\begin{eqnarray}
        x\left[ n\right] &=&Xz^n, \\
        y\left[ n\right] &=&Yz^n, \\
        y\left[ n+1\right] &=&Yz^{n+1}, \\
        y\left[ n+2\right] &=&Yz^{n+2}.
\end{eqnarray}
Reemplazamos en la ecuaci\'{o}n de diferencias 
\begin{equation}
        z^{2}Yz^n+azYz^n+bYz^n=cXz^n.
\end{equation}
Simplificamos y agrupamos 
\begin{equation}
        \left( z^{2}+az+b\right) Y=cX.
\end{equation}
Despejando $\frac{Y}{X}$ resulta 
\begin{equation}
        \frac{Y}{X}=H\left( z\right) =\frac{c}{\left( z^{2}+az+b\right) }.
\end{equation}
Observe que si $d_{1}\neq d_{2}$ (los polos de $H\left( z\right) )$ son distintos entre s\'{\i} entonces 
\begin{equation}
        H\left( s\right) =\frac{r_{1}}{z+d_{1}}+\frac{r_{2}}{z+d_{2}}.
\end{equation}
Las constantes $r_{1}$ y $r_{2}$ pueden calcularse 
\begin{eqnarray}
\left. H\left( z\right) \left( z+d_{1}\right) \right| _{z=-d_{1}} &=&
\left.\left( \frac{r_{1}}{z+d_{1}}+\frac{r_{2}}{z+d_{2}}\right) \left(z+d_{1}\right) \right| _{z=-d_{1}},\\
\left. H\left( z\right) \left( z+d_{2}\right) \right| _{z=-d_{2}} &=&
\left.\left( \frac{r_{1}}{z+d_{1}}+\frac{r_{2}}{z+d_{2}}\right) \left(z+d_{2}\right) \right| _{z=-d_{2}}.
\end{eqnarray}
Entonces podemos calcular la secuencia $y\left[ n\right] =h\left[ n\right] $ respuesta a una entrada
\begin{equation} 
        x\left[ n\right] =\delta \left[ n\right] =
                \left\{ \begin{array}{c}
                        n=0\rightarrow 1 \\ 
                        n\neq 0\rightarrow 0
                \end{array} \right.
\end{equation} como 
\begin{equation}
h\left[ n\right] =\left( r_{1}\left( -d_{1}\right) ^n+r_{2}\left(-d_{2}\right) ^n\right) u\left[ n\right].
\end{equation}

\begin{enumerate}
\item   Si $\left| d_{1}\right| <1$ y $\left| d_{2}\right| <1$ simult\'{a}neamente, entonces el sistema es 
estable porque
        \begin{equation}
                \begin{array}{c}
                        Lim \\ 
                        n\rightarrow \infty
                \end{array}
                h\left[ n\right] =0.
        \end{equation}
        
\item   En cambio, si $\left| d_{1}\right| >1$ o $\left|d_{2}\right| >1,$ entonces el sistema es inestable 
porque 
        \begin{equation}
                \begin{array}{c}
                        Lim \\ 
                        n\rightarrow \infty
                \end{array}
                h\left[ n\right] =\infty .
        \end{equation}
\end{enumerate}

\section{Aplicaci\'{o}n}

\subsection{Red el\'{e}ctrica en escalera}

En la figura \ref{laddercircuiteps}, $r$ es la resistencia serie y $g$
es la conductancia shunt. Aplicando la ley de Kirchoff de las
corrientes en el nodo central tenemos 
\begin{equation}
        \frac{v_{0}\left[ n-1\right] - v_{0}\left[ n\right] }{r}+
        \frac{v_{0}\left[ n+1\right] - v_{0}\left[ n\right] }{r}+
        g\left( v_i\left[ n\right] -v_o\left[n\right] \right) = 0,
\end{equation}
de donde se obtiene la ecuaci\'{o}n de diferencias
\begin{equation}
        -v_{0}\left[ n+1\right] +
        \left( 2+rg\right) v_{0}\left[ n\right] -
        v_{0}\left[n-1\right] =
        rgv_i\left[ n\right].  \label{sisCuatro}
\end{equation}

\begin{figure}
  \centering
  \includegraphics[height=2in]{./graphics/ladder_circuit.eps}
  \caption{Circuito escalera}
  \label{laddercircuiteps}
\end{figure}

Falta completar

